\documentclass{report}

\usepackage{amsmath}
\usepackage{amsthm}
\usepackage{amssymb}
\usepackage{mathpartir}

\title{Polymorphic bottom-up weighted relational programming}
\author{Dmitri Volkov}

\newcommand{\wt}{\mathrel{\text{wt}}}
\newcommand{\vt}{\mathrel{\text{vt}}}

\begin{document}

\maketitle

\chapter{Polymorphic semiringKanren}

\section{Typing rules}

We expand the syntax and typing of minimal semiringKanren to support polymorphic relations.

\begin{figure}
\[
\begin{array}{lrcl}
	\text{Programs} & P & ::= & \mathit{Rel}\dots \\
	\text{Relations} & Rel & ::= & \texttt{(defrel (\(R\) \(\forall \alpha \dots\) . (\(x:\tau\)) \(\dots\)) \(g\)) } \\
	\text{Goals} & g & ::= & \texttt{(conj \(g\) \(g\))} \\
	& & | & \texttt{(disj \(g\) \(g\))} \\
	& & | & \texttt{(fresh ((\(x:\tau\)) \(\dots\)) \(g\))} \\
	& & | & \texttt{(== \(v\) \(v\))} \\
	& & | & \texttt{(=/= \(v\) \(v\))} \\
	& & | & \texttt{(\(R\) \(v\) \(\dots\))} \\
	& & | & \texttt{(factor \(k\))} \\
	\text{Values} & v & ::= & \texttt{sole} \\
	& & | & \texttt{(left\(_{\tau_1, \tau_2}\) \(v\))} \\
	& & | & \texttt{(right\(_{\tau_1, \tau_2}\) \(v\))} \\
	& & | & \texttt{(pair \(v\) \(v\))} \\
	& & | & x \\
	\text{Types} & \tau & ::= & \texttt{Unit} \\
	& & | & \texttt{(Sum \(\tau\) \(\tau\))} \\
	& & | & \texttt{(Pair \(\tau\) \(\tau\))} \\
	& & | & \alpha \\
	\text{Relation names} & R & & \\
	\text{Variables} & x,y,z & & \\
	\text{Type variables} & \alpha,\beta,\gamma & & \\
	\text{Weights} & r & \in & \mathbb{K}
\end{array}
\]
\end{figure}

Here, the \(\Gamma\) context holds relation types, and the \(\Delta\) context holds variable types.
The \(\Delta \vdash v:\tau\) is the value (and therefore variable) typing judgement asserting that value \(v\) has type \(\tau\).
The \(\Gamma;\Delta \vdash g \wt\) judgement denotes that the goal \(g\) is well-typed under the \(\Gamma\) and \(\Delta\) contexts.
The \(\Gamma \vdash R \wt\) judgement denotes that the relation \(R\) is well-typed under the \(\Gamma\) context, and the \(P \wt\) judgement denotes that the entire program is well-typed.
The \(\Delta \vdash \tau \vt\) judgement denotes that the type \(\tau\) is valid under the \(\Delta\) context.
The judgements within the \(\Gamma\) context are extended to support type variables: \((R : \forall \alpha_1, \dots, \alpha_m . \tau_1,\dots,\tau_n \rightarrow)\).
As before, \(\Delta\) holds value typing judgements (of the form \(x:\tau\)). 
\(Delta\) is also extended to hold judgements of the form \(\alpha : *\), indicating that the type variable \(\alpha\) can be treated as a valid type in the current context.
In both cases, note that type variables may now appear in \(\tau\).

Note that our relation typing judgement is "the relation is well-typed under a context" rather than "the relation has type \(\dots\)".
Relations are required to specify all needed typing information, so it only is necessary to check that their bodies are compatible.
Because the types of all relations are known ahead-of-time, we avoid potential issues of recursive dependencies when type-checking relation calls.
Also note that in the formal semantics type variables are explicitly specified, but in practice they can be automatically extracted by looking at the argument types.

\begin{figure}
\begin{mathpar}
	\infer{ }{\Delta \vdash \texttt{Unit} \vt}
	\and
	\infer{
		\Delta \vdash \tau_1 \vt \and \Delta \vdash \tau_2 \vt
	}{
		\Delta \vdash \texttt{(Sum \(\tau_1\) \(\tau_2\))} \vt
	}
	\and
	\infer{
		\Delta \vdash \tau_1 \vt \and \Delta \vdash \tau_2 \vt
	}{
		\Delta \vdash \texttt{(Prod \(\tau_1\) \(\tau_2\))} \vt
	}
	\and
	\infer{ }{\Delta,\alpha : * \vdash \alpha \vt}
\end{mathpar}
\caption{Inference rules for eype validity judgement.}
\end{figure}

\begin{figure}
\begin{mathpar}
	\infer{ }{
		\Delta \vdash \texttt{sole} : \texttt{Unit}
	}
	\and
	\infer{
		\Delta \vdash v : \tau_1
	}
	{
		\Delta \vdash \texttt{(left\(_{\tau_1, \tau_2}\) \(v\))} : \texttt{(Sum \(\tau_1\) \(\tau_2\))}
	}
	\and
	\infer{
		\Delta \vdash v : \tau_2
	}
	{
		\Delta \vdash \texttt{(right\(_{\tau_1, \tau_2}\) \(v\))} : \texttt{(Sum \(\tau_1\) \(\tau_2\))}
	}
	\and
	\infer{
		\Delta \vdash v_1 : \tau_1
		\and
		\Delta \vdash v_2 : \tau_2
	}
	{
		\Delta \vdash \texttt{(pair \(v_1\) \(v_2\))} : \texttt{(Prod \(\tau_1\) \(\tau_2\))}
	}
	\and
	\infer{
		\Delta, x : \tau \vdash \tau \vt
	}
	{
		\Delta, x : \tau \vdash x : \tau
	}
\end{mathpar}
\caption{Inference rules for value typing judgement.}
\end{figure}

\begin{figure}
\begin{mathpar}
	\infer{
		\Gamma;\Delta \vdash g_1 \wt
		\and
		\Gamma;\Delta \vdash g_2 \wt
	}{
		\Gamma;\Delta \vdash \texttt{(conj \(g_1\) \(g_2\))} \wt
	}
	\and
	\infer{
		\Gamma;\Delta \vdash g_1 \wt
		\and
		\Gamma;\Delta \vdash g_2 \wt
	}{
		\Gamma;\Delta \vdash \texttt{(disj \(g_1\) \(g_2\))} \wt
	}
	\and
	\infer{
		\Delta \vdash \tau \vt
		\and
		\Gamma;\Delta,x:\tau \vdash g \wt
	}{
		\Gamma;\Delta \vdash \texttt{(fresh ((\(x:\tau\))) \(g\))} \wt
	}
	\and
	\infer{
		\Delta \vdash v_1 : \tau
		\and
		\Delta \vdash v_2 : \tau
	}{
		\Gamma; \Delta \vdash \texttt{(== \(v_1\) \(v_2\))} \wt
	}
	\and
	\infer{
		\Delta \vdash v_1 : \tau
		\and
		\Delta \vdash v_2 : \tau
	}{
		\Gamma; \Delta \vdash \texttt{(=/= \(v_1\) \(v_2\))} \wt
	}
	\and
	\infer{
	}{
		\Gamma, (R : \tau_1, \dots, \tau_n \rightarrow); \Delta, v_1 : \sigma(\tau_1), \dots, v_n : \sigma(\tau_n) \vdash \texttt{(\(R\) \(v_1\) \(\dots\) \(v_n\))} \wt
	}
	\and
	\infer{ }{
		\Gamma; \Delta \vdash \texttt{(factor \(r\))} \wt
	}
\end{mathpar}
\caption{Inference rules for goal typing judgement.}
\end{figure}

\begin{figure}
\begin{mathpar}
	\infer{
		\Gamma ; \alpha_1:*, \dots, \alpha_m:*, x_1:\tau_1, \dots, x_n : \tau_n \vdash g \wt
	}{
		\Gamma \vdash \texttt{(defrel (\(R\) \(\forall \alpha_1, \dots, \alpha_m\) . (\(x_1 : \tau_1\)) \(\dots\) (\(x_n : \tau_n\))) \(g\)))} \wt
	}
	\and
	\infer{
		(R_1 : \forall \alpha_1, \dots, \alpha_m . \tau_1, \dots \tau_n),
		\dots,
		(R_N : \forall \gamma_1, \dots, \gamma_q . \rho_1, \dots, \rho_p)
		\vdash \\\\
		\texttt{(defrel (\(R_1\) \(\forall \alpha_1, \dots, \alpha_m\) . (\(x_1 : \tau_1\)) \(\dots\) (\(x_n : \tau_n\))) \(g_1\)))} \wt
		\\\\
		\and
		\vdots
		\and
		\\\\
		(R_1 : \forall \alpha_1, \dots, \alpha_m . \tau_1, \dots \tau_n),
		\dots,
		(R_N : \forall \gamma_1, \dots, \gamma_q . \rho_1, \dots, \rho_p)
		\vdash \\\\
		\texttt{(defrel (\(R_N\) \(\forall \gamma_1, \dots, \gamma_q\) . (\(z_1 : \rho_1\)) \(\dots\) (\(z_p : \rho_p\))) \(g_N\)))} \wt
	}{
		\texttt{(defrel (\(R_1\) \(\forall \alpha_1, \dots, \alpha_m\) . (\(x_1 : \tau_1\)) \(\dots\) (\(x_n : \tau_n\))) \(g_1\)))}, \dots, \\\\
		\texttt{(defrel (\(R_N\) \(\forall \gamma_1, \dots, \gamma_q\) . (\(z_1 : \rho_1\)) \(\dots\) (\(z_p : \rho_p\))) \(g_N\)))} \wt \\\\
	}
\end{mathpar}
\caption{Inference rules for relation and program typing judgements.}
\end{figure}

\end{document}
